\subsection{Regularization, Normalization, and Generalization}

Deep-learning regularization combines penalties, stochastic training effects, architecture, and validation-based stopping.

\begin{definition} \label{def:sl-neural-regularized-objective} \hlt{Parameter-Norm Regularization}\\
Given objective $J(\theta)$, a regularized neural-network objective is
\[
\widetilde{J}(\theta)=J(\theta)+\alpha\Omega(\theta),
\qquad
\alpha\geq0.
\]
Weight decay uses an $\ell^2$ penalty on weights, while $\ell^1$ penalties encourage sparsity.
\end{definition}

\begin{definition} \label{def:sl-dropout} \hlt{Dropout}\\
Dropout trains subnetworks obtained by multiplying input or hidden units by random Bernoulli masks.  It can be interpreted as optimizing an average loss over a large tied-parameter ensemble.
\end{definition}

\begin{method} \label{method:sl-early-stopping} \hlt{Early Stopping}\\
During training, track validation loss and return the parameter value with the best validation performance.  Early stopping limits the reachable region of parameter space and can act like an implicit regularizer.
\end{method}

\begin{definition} \label{def:sl-batch-normalization} \hlt{Batch Normalization}\\
Batch normalization replaces a minibatch of activations $H$ by standardized activations
\[
H'=\frac{H-\mu}{\sigma},
\]
then often applies learned scale and shift parameters.  At evaluation time, running averages of $\mu$ and $\sigma$ replace minibatch statistics.
\end{definition}

\begin{remark} \label{rem:sl-data-augmentation-noise} \hlt{Augmentation and Noise}\\
Data augmentation, label smoothing, noise injection, adversarial training, and multitask learning regularize by changing the effective training distribution or by constraining learned representations.
\end{remark}
